\documentclass[conference]{IEEEtran}
\IEEEoverridecommandlockouts
\usepackage{cite}
\usepackage{amsmath,amssymb,amsfonts}
\usepackage{algorithmic}
\usepackage{graphicx}
\usepackage{textcomp}
\usepackage{xcolor}
\usepackage{hyperref}
\usepackage{float}
\usepackage{booktabs}

\def\BibTeX{{\rm B\kern-.05em{\sc i\kern-.025em b}\kern-.08em
    T\kern-.1667em\lower.7ex\hbox{E}\kern-.125emX}}

\begin{document}

\title{Rescue Rover: Autonomous Obstacle Detection and Environment Mapping using ESP32 and Ultrasonic Sensing}

\author{\IEEEauthorblockN{Anshaditya Sharma}
\IEEEauthorblockA{\textit{23BRS1204}\\
}
\and
\IEEEauthorblockN{Agrim Gusain}
\IEEEauthorblockA{\textit{23BRS1080}\\
}
}

\maketitle

\begin{abstract}
Traditional Search and Rescue (SAR) operations are severely limited by environments characterized by zero visibility and structural instability. While previous research has focused on modular controllers for drones and quadrupeds using optical 3D reconstruction (NeRFs), these systems fail in pitch-black or smoke-filled disaster zones. This report proposes an autonomous rescue rover based on the ESP32 microcontroller that utilizes ultrasonic sensor arrays for real-time obstacle detection and environment mapping. By measuring time-of-flight (ToF) acoustic pulses, the system builds spatial representations of surroundings where human reach and light-based sensors are non-functional. The methodology involves fusing data from multiple HC-SR04 sensors to map a 2D proximity grid and executing local path planning logic on the edge. The integration creates a highly scalable, low-cost platform that functions precisely in critical dark-zone environments.
\end{abstract}

\begin{IEEEkeywords}
Search and Rescue, ESP32, Ultrasonic Sensing, Autonomous Navigation, Obstacle Detection, Edge Computing
\end{IEEEkeywords}

\section{Introduction}
As established in recent SAR research, the first 72 hours of a disaster are historically recognized as the ``golden window'' for survivor recovery. However, traditional human responders and even advanced visually-oriented robots face significant hurdles in ``dark-zone'' operations—areas such as collapsed basements, tunnels, or smoke-filled corridors where ambient light is entirely absent and physical human entry is too perilous. The sheer complexities of post-disaster environments demand highly durable, versatile, and independent robotic systems to safely execute mapping and victim localization duties.

While concurrent research in this domain introduces a ``Universal Platform'' for modular control of heterogeneous robots, our proposed system extends this fundamental logic to the embedded hardware level through the ESP32 microcontroller architecture. Rather than relying on computationally heavy, light-dependent optical Neural Radiance Fields (NeRFs) or visual Simultaneous Localization and Mapping (vSLAM), we pivot toward ultrasonic ToF (Time-of-Flight) based distance measurement. This modality enables resilient spatial awareness through acoustic echo patterns, ensuring that neither total darkness nor debris-induced visual occlusion prevents accurate mapping of the operation theatre.

\section{System Architecture}
The rover is deliberately designed as a high-resilience, low-cost micro-node within a broader, coordinated SAR network. The decentralized nature of its computing ensures that in the event a rover is incapacitated, the global mapping model is minimally affected.

The fundamental components of the architecture include:
\begin{itemize}
    \item \textbf{Processing Hub:} The ESP32 module (WROOM-32) delivers dual-core processing essential for the rapid concurrent execution of real-time signal acquisition, data filtering, and WiFi transmission.
    \item \textbf{Sensing Modality:} Instead of fragile RGB stereoscopic cameras or costly LiDAR modules, the system arrays multiple HC-SR04 ultrasonic logic units. These sensors calculate ToF acoustic reflections, mapping physical obstacles dynamically.
    \item \textbf{Kinematics \& Motion:} An L298N dual H-Bridge motor driver receives High-Frequency Pulse Width Modulated (PWM) signals from the ESP32, translating directional logic into 4-wheel drive articulation.
    \item \textbf{Communication Layer:} The built-in 802.11 b/g/n WiFi transceiver on the ESP32 establishes a telemetry bridge to a ground station over WebSocket or UDP protocols. 
\end{itemize}

\begin{figure}[H]
    \centering
    \includegraphics[width=\linewidth]{architecture.png}
    \caption{Rescue Rover System Architecture — Four-layer data flow with base station telemetry communication.}
    \label{fig:architecture}
\end{figure}

\section{Hardware Implementation}

To construct an autonomous rover capable of enduring rubble traversal, a sturdy acrylic or 3D-printed PLA base houses the internal electronics securely. The power delivery system relies on a dual-cell 7.4V Lithium-Ion package to guarantee substantial continuous discharge for the DC gear motors while preserving stable logical voltages.

\begin{figure}[H]
    \centering
    \includegraphics[width=\linewidth]{circuit.png}
    \caption{Rescue Rover Circuit and Electronic Block Diagram.}
    \label{fig:circuit}
\end{figure}

\begin{table}[H]
\centering
\caption{System Hardware Specifications}
\renewcommand{\arraystretch}{1.5}
\begin{tabular}{|l|l|}
\hline
\textbf{Component Level} & \textbf{Specification / Detail} \\
\hline
\textbf{Microcontroller} & ESP32 Dual-Core @ 240MHz, 520KB SRAM \\
\hline
\textbf{Acoustic Sensing} & 3x HC-SR04 (40kHz freq, 2cm-400cm range) \\
\hline
\textbf{Motion Actuation} & 4x 6V DC Gear Motors (BO Motors) \\
\hline
\textbf{Orientation/IMU} & MPU6050 (6-Axis Accelerometer + Gyro) \\
\hline
\textbf{Motor Driver} & L298N Dual H-Bridge (2A peak per channel) \\
\hline
\textbf{Power Supply} & 2x 18650 3.7V Cells in Series (7.4V total) \\
\hline
\end{tabular}
\label{tab:hardware_specs}
\end{table}

\begin{figure}[H]
    \centering
    \includegraphics[width=\linewidth]{rover_car_gen.png}
    \caption{Rescue Rover prototype with prominent ultrasonic sensors navigating a simulated dark disaster zone environment.}
    \label{fig:rover_gen}
\end{figure}

\section{Methodology}
\subsection{Ultrasonic Time-of-Flight Dynamics}
Unlike optical sensors that detect photon reflections, acoustic proximity mapping relies on piezoelectric transducers emitting 40kHz ultrasound bursts. When encountering an object, the acoustic wave reflects and propagates back to the receiver. The distance $d$ to an obstacle is governed by the Time-of-Flight ($t$) and the speed of sound in air ($v_{sound} \approx 343$ m/s at 20$^\circ$C):

\begin{equation}
    d = \frac{v_{sound} \times t}{2}
\end{equation}

where the division by two accounts for the round-trip travel time of the ultrasonic pulse. By positioning sensors at the forward, left, and right chassis orientations, a minimal grid matrix of immediate physical blockages can be continuously refreshed at ~20Hz intervals.

\subsection{Path Planning and Autonomous Navigation}
Distance constraints from the ultrasonic tri-array feed into a computationally conservative reactive state machine executing onboard the ESP32. The algorithmic loop follows a hierarchical logic structure:
\begin{enumerate}
    \item \textbf{Collision Imminence Detection:} Continuously polling the forward HC-SR04. If the distance falls directly below the hyper-parameter threshold ($d < 20$ cm), emergency braking is immediately engaged to protect hardware integrity.
    \item \textbf{Directional Arbitration:} Once stopped, the system assesses the vectors calculated by the lateral (Left/Right) sensors. The microcontroller parses the arrays to find the direction with the maximum unobstructed run-length:
    $$ D_{next} = \arg\max (d_{left}, d_{right}) $$
    \item \textbf{Motor Actuation Output:} Command sequences map to the L298N driver, reversing motor parity on specific sides to execute an in-place zero-radius turn toward the calculated $D_{next}$.
    \item \textbf{Grid Transmittal:} Simultaneously, distance nodes and relative orientation vectors (calculated via the integrated MPU6050 IMU using a complementary filter) are published to the network for remote multi-dimensional mapping.
\end{enumerate}

\section{System Analysis and Validation}
The utility of the sonar-centric methodology is clearest when benchmarked against parallel optic-based Search and Rescue robotic units. The fundamental performance paradigm dictates that different sensors handle distinct subsets of disaster phenomenology better than others. 

\begin{table}[H]
\centering
\caption{Comparison with Alternative SAR Robotic Approaches}
\renewcommand{\arraystretch}{1.5}
\begin{tabular}{|p{2.3cm}|p{2.6cm}|p{2.7cm}|}
\hline
\textbf{Assessment Metric} & \textbf{Visual/NeRF Platform (Base Paper)} & \textbf{Ultrasonic Platform (Proposed ESP32 Node)} \\
\hline
\textbf{Visibility Needs} & High (Extensive ambient light necessary) & Zero (Acoustic operation independent of lux) \\
\hline
\textbf{Hardware/Compute} & High (Requires GPU-accelerated backend) & Minimal (Calculations executed on ESP32 Edge) \\
\hline
\textbf{Interference Resist} & Susceptible to smoke, excessive dust, steam & High resilience to airborne dust/smoke particulate\\
\hline
\textbf{Production Cost} & Extremely High (Boston Dynamics Spot / DJI) & Sub-\$50 (Mass deployable logic boards) \\
\hline
\textbf{Field Application} & Broad aerial reconstruction / High altitude & Narrow, localized inner-structure traversal \\
\hline
\end{tabular}
\label{tab:comparison}
\end{table}

\begin{figure}[H]
    \centering
    \includegraphics[width=\linewidth]{comparison.png}
    \caption{Visual slide presentation format of the ESP32 Ultrasonic Platform vs traditional visual processing platforms.}
    \label{fig:comparison_slide}
\end{figure}

\section{Challenges and Potentials}
Operating within the unpredictable terrain of disaster wreckage inherently introduces distinct acoustic challenges. Extremely highly porous, soft materials (e.g., loose insulation fabric, mattresses) may absorb the 40kHz signal, resulting in weak echoes that the controller might falsely interpret as empty space. Additionally, angled, highly smooth surfaces might deflect sound waves away from the receiving transducer (specular reflection), creating false negatives in mapping matrices.

However, resolving these issues lies in the \textit{system-of-systems} approach. Rather than acting as an isolated navigator, the ESP32 rover effectively acts as a ``proximate sensory extension cable,'' uploading sparse, localized, zero-visibility acoustic data to the centralized SAR network. Ground stations can probabilistically combine this data with overarching visual models. Should this acoustic layer be fused with the visual Universal Platform, an all-encompassing multi-modal map of the collapsed structure is generated, allowing teams to maneuver in the most treacherous regions.

\section{Conclusion}
The integration of ESP32-mediated ultrasonic rangefinding with decentralized path-planning logic introduces a robust paradigm shift for Search and Rescue robotics specifically operating in zero-visibility scenarios. By eschewing computationally costly and light-reliant visual processing mechanisms, the Rescue Rover maintains critical navigational capabilities precisely where complex algorithms fail. Given the microscopic manufacturing costs and formidable mapping precision of acoustic time-of-flight technology, these swarms of ESP32 edge nodes are poised to dramatically supplement wider disaster recovery networks, safely mapping the unknowns and expediting survivor extraction against critical timelines.

\begin{thebibliography}{00}
\bibitem{b1} D. Cafolla, B.D.M. Chaparro-Rico, X. Xie, ``Design and testing of a universal platform for search and rescue operation: Exploring indoor and outdoor potentials,'' \textit{Robotics and Autonomous Systems}, vol. 196, 2026, Article No. 105237.
\bibitem{b2} J. Borenstein, Y. Koren, ``The vector field histogram-fast obstacle avoidance for mobile robots,'' \textit{IEEE Transactions on Robotics and Automation}, vol. 7, no. 3, pp. 278-288, June 1991.
\end{thebibliography}

\end{document}
